\begin{ans}{1.2.1}
  略.
\end{ans}

\begin{ans}{1.2.2}
  略.
\end{ans}

\begin{ans}{1.2.3}
  略.
\end{ans}

\begin{ans}{1.2.4}
  (合成数の定義は正の数に限る場合があるが, ここでは負の数も含むことにする.)
  十分大きい自然数$m$に対して$\abs{f(m)} > 1$となる. このような$m$を用いて
  \[
    f(x) = q(x)(x - m) + r
  \]
  とおくと, $\abs{r} = \abs{f(m)} > 1$である.
  $n (> 0)$として, ある正の整数$k$を用いて$n = k\abs{r} + m$と書けて, かつ$f(x)$, $f(x)-r$, $f(x)+r$のいずれの根でもないもの (これらの多項式は$0$ではなく, 根の個数はたかだか有限なので存在する) をとれば,
  $f(n)$は$\abs{r} (> 1)$の倍数かつ$0$, $\pm r$でないので合成数である.
\end{ans}

\begin{ans}{1.2.5}
  $M = \{ f \mid f: I \rightarrow \mathbb{N}, \text{有限個の } i \in I \text{を除き } f(i) = 0 \}$
  とおくと, 値による和をモノイド演算として$M$はモノイドとなる.
  このモノイド$M$に対するモノイド環$A[M]$が, $I$を添え字集合とする無限変数多項式環にほかならない.
\end{ans}
